\section{Формальная постановка задачи}
\label{sec:2_problem_statement} \index{2_problem_statement}

\subsection{Модель прикладной программы}
Модель прикладной программы задается ориентированным ациклическим графом потока данных
\[
G = (V, E), \quad |V| = n, \quad |E| = m.
\]
Каждой вершине $v \in V$ соответствует работа, а каждой дуге $(u, v) \in E$ — зависимость по данным между работами $u$ и $v$. Если $(u, v) \in E$, то работа $v$ может начать выполняться только после завершения работы $u$.

Для каждой работы $v \in V$ задана длительность исполнения $d_v > 0$. Кроме того, каждая работа может формировать один или несколько выходных буферов. Обозначим множество буферов, создаваемых работой $v$, через
\[
B(v) = \{b_{v1}, b_{v2}, \dots, b_{vk}\}.
\]
Для каждого буфера $b \in B(v)$ задаются его объем памяти $r_b > 0$ и множество потребителей $C(b) \subseteq V$. Если $u \in C(b)$, то работа $u$ использует результат, содержащийся в буфере $b$.

\subsection{Модель вычислительной системы}
Вычислительная система состоит из $P$ одинаковых процессоров
\[
SP = \{sp_1, sp_2, \dots, sp_P\}.
\]
Каждый процессор в любой момент времени может выполнять не более одной работы. Прерывания работ не допускаются, миграция уже запущенной работы между процессорами невозможна.

Задержки межпроцессорной передачи данных не учитываются. Поэтому момент готовности работы к выполнению определяется завершением всех ее непосредственных предков. Ограничивающими факторами являются зависимости между работами, число доступных процессоров и общий объем памяти $M$, доступный для хранения активных буферов.

\subsection{Модель использования памяти}
Пусть работа $v$ стартует в момент $s_v$ и завершается в момент
\[
f_v = s_v + d_v.
\]
Буфер $b \in B(v)$ выделяется в момент старта работы $v$. В данной модели это означает предварительное резервирование памяти под результат работы: перед запуском работы система должна убедиться, что памяти достаточно для ее выходных буферов, и только после этого работа может быть начата. Освобождение буфера происходит после завершения последней из работ-потребителей его содержимого (в их вершины идут дуги, помеченные этим буфером):
\[
fr_b = \max_{u \in C(b)} f_u.
\]

Обозначим через $J(\tau)$ множество буферов, активных в момент времени $\tau$. Тогда текущее использование памяти равно
\[
\operatorname{mem}(\tau) = \sum_{b \in J(\tau)} r_b.
\]
Пиковое использование памяти на расписании $HP$ определяется как
\[
\operatorname{peak}(HP) = \max_{0 \le \tau \le makespan} (\operatorname{mem}(\tau)).
\]
Расписание считается допустимым по памяти, если выполнено условие
\[
\operatorname{peak}(HP) \le M.
\]

\subsection{Модель расписания}
Расписание $HP$ задает для каждой работы $v \in V$:
\begin{enumerate}
	\item номер процессора $p(v) \in \{1, \dots, P\}$;
	\item момент старта $s_v$;
	\item момент завершения $f_v = s_v + d_v$.
\end{enumerate}
Расписание является корректным, если выполняются следующие ограничения:
\begin{enumerate}
	\item каждая работа назначена ровно одному процессору;
	\item на каждом процессоре интервалы выполнения назначенных ему работ не пересекаются;
	\item для каждой дуги $(u, v) \in E$ выполнено неравенство $f_u \le s_v$;
	\item для всего расписания выполнено ограничение по памяти $\operatorname{peak}(HP) \le M$.
\end{enumerate}

\subsection{Минимизируемая целевая функция}
Минимизируемым критерием является длительность расписания, то есть момент завершения последней работы:
\[
C_{\max}(HP) = \max_{v \in V} f_v.
\]
Требуется найти такое допустимое расписание $HP^*$, что
\[
C_{\max}(HP^*) = \min_{HP \in \mathcal{H}_{\text{adm}}} C_{\max}(HP),
\]
где $\mathcal{H}_{\text{adm}}$ — множество всех корректных расписаний.

Похожие задачи рассматривались в работах~\cite{Wang2010},~\cite{Herrmann2014} и~\cite{Venugopalan2012}, а в работе~\cite{Papp2025} доказывается, что подобная задача является NP-трудной.